\chapter{Introduction}

Hello, and welcome to the \textbf{Documentary of Modelling Automata}. The name might already suggest something to you, but I will leave that toward your own interpretation on what we do, akin to a \textit{Revue D'études Scientifiques}, sort of. 

Generally, this book serves a dual purpose - one is to record my understanding and knowledge about the topic, and related to it, and one is to develop my own theories and models - most likely original, more often than not it's not. Also, it will also include rough implementation, papers' ideas analysis, and else. So, a bag full of everything, I suppose.  of a twofold reason - one is to record my understanding and knowledge about the topic, and related to it, and one to directly record and write down my discovery, or rather, most of the time just analysis and designs. 

It might sound as if this book is entirely self-sufficient, and also self-serving, yet it isn't. The text cannot take into account all preliminaries or requirements of understanding, yet. It is simply too large to take all in, because the field itself is very complex, and perhaps also of a more ergonomic issue since it will also present new results or theory in companion. On the side of the narrative in favour, perhaps it can be broken down to $0-1$ loss function too (it's a joke), since it is perhaps too difficult to follow directly from the get-go, so I pretend to be the third-party narrative. With that, the book is not just about myself asking and answering questions, logging knowledge, but to also write and try to explain it as a perhaps to a 5-year-old, to somehow make it sounds not like listening to quantum mechanics. 

Overall, I seriously hope that this book will help me myself, ultimately, and help anyone along the way, if they stumble upon this. Hopefully you will give me some exposure too (famous?), so, I am counting on you. With that said, artificial intelligence is a hoax. Only us can refute that. 

\section{The motivation of existence}

The motivation of how this get here in the first place is rather simple. Perhaps more straightforward than you would imagine. It started simple as a book to comprise knowledge, in the sense that it debates between the philosophy of knowledge centralization - books and materials, linearity organization (not in the sense of mathematics), and such forth, and the philosophy of \textit{zettelkasten}, or by the little notes and box, decentralization and of connections, dynamic and fluid yet disorganized in the sense of a pathway. Ultimately, I chose to do both, and the byproduct of such judgement is this and a few more precursors of it. What kind of knowledge are there do you ask? Well, at first, nothing too crazy. Mostly just thoughts and some topics that I think I should know of, like mathematics and the like. Soon it turns into a place where I also put my own inquiry of mathematics on there, and philosophy of such. Actually, there is an article about philosophy of probability too, which is fairly erratic considering that I learnt about probability theory 4 weeks afterward of that article. 

Soon enough, the scope it wider and wider. Until it becomes the first precursor of this book as of today, the content went by a dramatic shift into the world of artificial intelligence, where we discuss and formulate the structure of which is then to be said exhibit the behaviours of the intelligent, yet artificially made. Such inquiry indeed took a lot of time to formulate and consider, and hence, renditions of the book exist in various sense, yet none seems particularly finished. Which is to say that I did not finish any of them. Some of it indeed led to questions and ideas that is then adopted to this book, explanation of certain problems that is fairly critical and else, but they indeed did not help much otherwise in creating a centralized theory just as I hoped it too. I mean, it also did give me a \LaTeX template to work on, so I guess it is not that bad, but the main purpose is undone. If I ever publish it, the subject would look more like a mess of convoluted pseudoscience mess, rather than a coherent picture. 

The main shift happened afterward, however, after that particular deadlock of over half a year to focus on something else entirely. It is at this time that I started to think more deeply, or more so broadly, to the problem of the construct, of the process, of the machines and its constructions, of the automata and its function, of the model and its encoding, of the theory and language that one creates and frame it upon such, and so much and more. All of which, resides in this book, or at least I hope so. Such theory broaden the generality of the concept of artificial intelligence, the modelling, and the mathematics. That is to say, the foundation and whatnot required to comprehend them, is very large. Those are the topics that I want to introduce, and unfortunately, the structure of the old book does not allow me to do such. Hence, this book as it currently stands. 

What I do is not simply original research, nor I want to claim that I make an entirely new claims, entirely new theory or of such. Maybe, but it would be damn well that my discovery and analysis is left in someone else's basement already, with far more time and thoughts put into such. What I want to do, is a radical and comprehensive view, insights, analysis from different perspective, picking up what has been left out of unrecognized of established Tower of Babel, and proceed from then upward. Such can only be done, through the looking glass, and with grace in mind, or the messiness of a feverish dream comes forth. Poetic as it is, boorish in reality shall it be realized, should I warn you to not take my words as something out of admiration of one's self work, but rather a cry on how much caffeine intake it took from the body to procure any of this. 

\section{What is in the book}
The book is concerned of the main umbrella topic of \textbf{artificial intelligences}. Alongside with it are relevant theories and practical implementations that support said theoretical view. So, you would also expect to have \textit{machine learning, mathematical modelling, a lot of mathematics, information theory, complexity analysis}, and more, for example, some exotic topics like genetic engineering principles, programming \textit{synthesis}, or something like temporal analysis of models. Aside from such, there is also neuroscience and other typical brain-related topics, which is fairly straightforward when talking about inclusion (not DEI more or less). Specifically, there will also be an entire large chapter on the formalism of the learning theory, in a rigorous sense, so many of the chapters would be there to reinforce this. One of the major topics in this book will also be the abstract \textbf{philosophical analysis} or in general, philosophy of subjects of interest, and topics on neuroscience and cognitive science. While it is not intended to be a comprehensive content, such chapters regarding of said content will be, and should be, inferred to certain sufficiency and of problems or dilemma closer to our own purposes, which will be, ironically, settled in a philosophical chapter by its own. Because of its purpose to not be a textbook, and closer to a self-thought book, the content will be filled with fillers or thoughts, sometimes rambling or long derivation or justification for examples and thereof. 

As it currently stands, the main top category is the parts, albeit presented implicitly. Comprehensively, there are three main parts. Part I on \textit{theory} - the supporting theory, discussions, results and analysis. Part II on \textit{advanced theory}, being called advanced just because it is my own implementation and theory in accordance. And part III on \textit{implementation} and any necessary details on such - so, it can contain sections on the Python language itself (which is boring) - but also sections concerning deployment or rather typical construction of what has been established in the preceding part. Though there are plans to go for C++ implementation, high-intensity computation is perhaps not in the list for the current time (April 2024-August 2025). This will have to couple with computation theory and the like to dissect, though some language dialect with type control would certainly help in both way, but require longer time to develop. For historical purpose, which will be the old implementation of precursor AI, some PROLOG would be expected. 

The style of this book is also fairly erratic, so to speak. Because of the huge debt toward old ideas and forgotten one in practice, the style of writing would focus on a lot of the historical development and motivations thereof, but also includes technical details and a fairly bit of hand-waiving sections, for some sections where prerequisite would take up more than the content it is planned to present. 
\section{Frequently (maybe) asked questions}
There are always inquiry questions to answer, and those to write of about itself. 

\subsection{How much topic is covered, generally?}
Well, not so much, but I hope it will be enough to formalize and construct a formal treatment of a potentially modern approach to artificial intelligence theory. Of the latest revision, however (May 18, 2025), the book has been subsequently changed of its status as a comprehensive AI-based book, to a rather general book by its own standard. Though, it does not contrive the structure for literally everything, by that we will narrow down to either philosophy, computer science theory, artificial intelligence, mathematics, and physics, with expansions to other topics beside the core interpretation.
\subsection{What's with the question section?}

The question section is rather weird - not so unique, but definitely abnormal comparing to normal book sequence. Here, the question section is not for simply `homework-style' reaffirming questions, but rather usually so set against such notion and more on inquiry questions, open questions, analytic questions which can have solutions to them or not, and many more. So expect some of those questions to be very implicit, sometimes not so well-worded, and sometimes very hard to take care of because of the nature of said questions. 

Nevertheless, there will be questions related to variations and derivations of different structures, models, in sections concerned of them. Most of them will be perhaps used throughout the text further on, or rather being development on proof-of-concept that can be advanced henceforth. I don't think we need answer key for a lot of them, since some are open questions, but I will provide what is available. 

\section{Construction}
Because the book is incomplete, this section then must be added for organizational purpose. And yes, without it, I would soon forget what to add and what to write into the manuscript itself, so that is for sure an indictment that I conjure myself upon here. Nevertheless, we need it for our template. 

\subsection{Topics in consideration}
We are considering adding \textit{neuroscience}: neuroanatomy up to neurophysiology, and computational neuroscience simulations. \textit{Advanced mathematical modelling principles}, plus some philosophical rework of the mathematical modelling scheme. \textit{Representation and encoding}, as well as the theory of language. We also miss the theory of automata, which is harder than it is. \textit{Novel learning theory} is in construction, but will have to rework for consideration. \textit{Philosophy and operations of computation} is also considered, as well as subsections on philosophy of abstractions. The general philosophical view is also lacking. 

Connections between sections and chapters are also in construction, as there exist no fundamental way to connect them altogether as for now. Organization at part-level also need consideration for the roadmap of the book, since a lot of contents overlap with each other. More so, the connection between the \textit{mathematics section} and the \textit{theory section} will have to be somewhat clear out. This is troublesome but doable to complete. 

Aside from that, we have (September 2025) a proposal for adding \textit{programming theory}, simulation theory (right before mathematical modelling), chaos theory, percolation theory, \textit{programming synthesis}, type theory (which goes after language), and variations of information theory, system theory, and philosophy of mathematics (in the philosophical section, though I suspect this is \textbf{the most important chapter} of the book due to its content). 
\subsection{Topics in working}
For now, they are all in the table of content, however, some hidden topic like \textit{automata theory} requires prerequisite in the theory of computation first. 